\section{Reproduction}
\label[appendix]{app:reproduction}

Everything in this monograph can be rebuilt from public sources with the project's code. The code is available at \url{https://github.com/maxime2476/Atlas} and archived on Zenodo (\href{https://doi.org/10.5281/zenodo.23305095}{doi:10.5281/zenodo.23305095}). This appendix says how, and what is needed.

\paragraph{Environments.} Python~3.12 is managed with \texttt{uv}. Its lock file fixes every package, resolved with no release later than \VerUvDate{} (for instance pandas~\VerPandas, scikit-learn~\VerSklearn, DuckDB~\VerDuckdb, GeoPandas~\VerGeopandas). R~\VerR{} is managed with \texttt{renv}, from a dated snapshot of the Posit package repository (\VerRepoDate), with \texttt{fixest}~\VerFixest. The two languages exchange data through Parquet files with documented schemas.

\paragraph{Raw data.} About \RawSizeGb~GB of raw files, never modified, of which \DriasSizeGb~GB of climate projections (NDVI, E-OBS and DRIAS included). A manifest lists \ManifestFiles{} files with their URL, access date and SHA-256 hash; an automated test checks every hash. The raw data are not redistributed with the code; the download script rebuilds them.

\paragraph{Pipeline.} The repository's \texttt{data/README.md} gives the commands in order:
\begin{enumerate}
  \item downloads and checks;
  \item weather, crop masks, exposure windows and panels;
  \item econometrics in R;
  \item forecasting and NDVI;
  \item the pre-specified analyses of phase 4 bis;
  \item the checks added later, including soils, irrigation volumes and climate projections;
  \item figures, then the tables and figures of this monograph.
\end{enumerate}
The longest steps are the NDVI processing (orders depend on the NASA queue, then about 15 minutes of aggregation and 30 minutes of forecasting) and the daily weather (about 20 minutes). Two sources require a free personal account: the NDVI (NASA Earthdata) and the climate projections (DRIAS and DRIAS-Eau, one account for both). DRIAS has no download interface: orders are placed through its web form, by a script that fills it in as a user would, and are delivered later as archives. Its terms forbid sharing the account, so a reproduction needs its own; the projections are therefore the one step that cannot be rebuilt without registration. All other sources are open without registration.

\paragraph{A clean rebuild.} Before the checks added later, the whole chain was run again on a fresh copy of the repository, from the raw files and without downloading anything. It gave the same intermediate data, the same result tables byte for byte (except the forecast tables, which differ by at most $5\times 10^{-15}$, floating-point rounding) and the same numbers as the text at that point. The analyses added afterwards (all those marked as added later in \cref{sec:maize,sec:forecasting}, and the projections of \cref{sec:discussion}) were run in the working repository, not on a fresh copy.

\paragraph{The numbers in this text.} Every estimate, $p$-value, test size and table in this monograph is generated: a script reads the result tables of the project and writes them as \LaTeX{} macros and tables. A few numbers are written in the text by hand: counts that define the design (17 and 83 departments, the numbers of years and clusters, the 300 placebo repetitions) and a few historical figures from the decision log, which are cited as such. The script also checks that the specifications still match their recorded hashes and that the valuation of 2003 is reproduced exactly; it stops otherwise.

\paragraph{Tests.} The test suite has \TestCount{} tests: \TestIntegrity{} integrity checks of the raw files, one per file, and \TestLogic{} tests of logic and data. Together they check schemas, keys, the cardinality of joins, value ranges, the degree-day formula, the absence of forecasting leakage, the agreement of the maize mask with official statistics, and the integrity of the raw files. Tests that need data are skipped when the data are absent. The R code has no unit tests; it is validated by comparisons with \texttt{fixest} and by the size simulations of \cref{sec:inference}.

\paragraph{Randomness.} All random draws use fixed seeds, recorded in the code. One departure from the seed written in the phase 4 bis specification is listed in \cref{app:deviations}.
